\documentclass[12pt,a4paper]{article}
% Drake_preamble.tex - LaTeX Preamble Template
% 作者: 謝慕揚, MD, PhD, FESC
% 
% 使用說明:
% 1. 表格請使用 longtable，不要有垂直分隔線 |
% 2. 表格內換行使用 \newline，不要使用 \\
% 3. 藥物名稱、檢驗、檢查請維持英文
% 4. Reference 格式請使用 \href 加上驗證過的 DOI

% Including essential packages for Chinese typesetting
\usepackage{xeCJK}
\usepackage{fontspec}
% Configuring Chinese font with Noto Serif CJK TC, fallback to SimSun
\IfFontExistsTF{Noto Serif CJK TC}{
  \setCJKmainfont{Noto Serif CJK TC}
  \setCJKsansfont{Noto Serif CJK TC}
  \setCJKmonofont{Noto Serif CJK TC}
}{\setCJKmainfont{SimSun}
  \setCJKsansfont{SimSun}
  \setCJKmonofont{SimSun}
}

% Including other necessary packages
\usepackage{geometry}
\usepackage{titlesec}
\usepackage{enumitem}
\usepackage{xcolor}
\usepackage{colortbl}
\usepackage{tcolorbox}
\usepackage{booktabs}
\usepackage{longtable}
\usepackage{array}
\usepackage{graphicx}
\usepackage{tikz}
\usepackage{fancyhdr}
\usepackage{multicol}
\usepackage{datetime2}
\usepackage{hyperref}
\usepackage{mdframed}
\usepackage{amsmath}
\usepackage{amssymb}
\usepackage{multirow}

\hypersetup{
    colorlinks=true,
    linkcolor=emergency_blue,
    citecolor=emergency_blue,
    urlcolor=emergency_blue,
    pdfborder={0 0 0}
}

% Defining page geometry
\geometry{
  top=2.5cm,
  bottom=2.5cm,
  left=2cm,
  right=2cm,
  headheight=25pt
}

% Adjusting topmargin to compensate for increased headheight
\addtolength{\topmargin}{-10pt}

% Defining colors
\definecolor{emergency_red}{RGB}{186,24,27}
\definecolor{emergency_blue}{RGB}{0,114,188}
\definecolor{emergency_gray}{RGB}{120,120,120}
\definecolor{highlight_yellow}{RGB}{255,250,205}

% Setting title formats
\titleformat{\section}[block]{\Large\bfseries\color{emergency_red}}{\thesection}{10pt}{}
\titleformat{\subsection}[block]{\large\bfseries\color{emergency_blue}}{\thesubsection}{8pt}{}
\titleformat{\subsubsection}[block]{\normalsize\bfseries\color{emergency_gray}}{\thesubsubsection}{6pt}{}

% Defining custom environment for important notes
\newmdenv[
    linecolor=emergency_red,
    backgroundcolor=highlight_yellow,
    linewidth=2pt,
    topline=true,
    bottomline=true,
    leftline=true,
    rightline=true,
    innertopmargin=10pt,
    innerbottommargin=10pt,
    innerleftmargin=10pt,
    innerrightmargin=10pt
]{importantbox}

% Defining note command with tcolorbox
\newcommand{\note}[1]{
    \par
    \noindent
    \begin{tcolorbox}[
        colback=emergency_blue!5!white,
        colframe=emergency_blue,
        title=\textbf{重點提醒},
        fonttitle=\bfseries,
        boxrule=0.5pt,
        arc=3pt,
        left=6pt,
        right=6pt,
        top=6pt,
        bottom=6pt
    ]
    #1
    \end{tcolorbox}
}

% Key findings box
\newcommand{\keyfinding}[1]{
    \par
    \noindent
    \begin{tcolorbox}[
        colback=yellow!10!white,
        colframe=orange!75!black,
        title=\textbf{核心發現摘要},
        fonttitle=\bfseries,
        boxrule=0.5pt,
        arc=3pt
    ]
    #1
    \end{tcolorbox}
}

% Defining custom date format for ISO 8601
\DTMnewdatestyle{iso}{
  \renewcommand{\DTMdisplaydate}[4]{\number##1-\number##2-\number##3}
  \renewcommand{\DTMDisplaydate}[4]{\number##1-\number##2-\number##3}
}
\DTMsetdatestyle{iso}
\newcommand{\mydate}{\DTMtoday}


% Custom header for this document
\pagestyle{fancy}
\fancyhf{}
\fancyhead[L]{\textcolor{emergency_red}{Intensive Care Medicine 教學講義}}
\fancyhead[R]{\textcolor{emergency_blue}{Diaphragm Ultrasound}}
\fancyfoot[C]{\textcolor{emergency_gray}{\thepage}}
\renewcommand{\headrulewidth}{1pt}
\renewcommand{\headrule}{\hbox to\headwidth{\color{emergency_red}\leaders\hrule height \headrulewidth\hfill}}

\begin{document}

\begin{center}
{\LARGE\bfseries\textcolor{emergency_red}{Intensive Care Medicine Editorial}}\\[0.5cm]
{\Large\textbf{How I Perform Diaphragmatic Ultrasound in the Intensive Care Unit}}\\[0.3cm]
{\large 如何在加護病房執行橫膈膜超音波}\\[0.5cm]
{\normalsize Intensive Care Med 2024;50:2175-2178. DOI: 10.1007/s00134-024-07688-x}\\[0.3cm]
{\normalsize 整理：謝慕揚 MD, PhD, FESC \quad 日期：\mydate}
\end{center}

\keyfinding{
\textbf{核心訊息：}Diaphragm ultrasound 是非侵入性床邊工具，特別適用於 difficult-to-wean 病人的診斷評估。

\textbf{實用重點：}
\begin{itemize}
    \item \textbf{Subcostal approach}：測量 diaphragm excursion (DE)，DE <10-15 mm 提示 dysfunction
    \item \textbf{Intercostal approach}：測量 thickness 及 thickening fraction (DTF)，DTF$_{max}$ <20\% 提示 weakness
    \item 評估時需移除或減少呼吸器支持，建議至少 40 次訓練才能獨立操作
\end{itemize}
}

\tableofcontents
\newpage

%==============================================================
\section{文獻概述}
%==============================================================

本篇為 Intensive Care Medicine 的 Editorial，由比利時 KU Leuven 的 Hermans 教授等人撰寫，分享 diaphragm ultrasound 在 ICU 的實務經驗與臨床應用。

\subsection{作者資訊}
\textbf{通訊作者：}Greet Hermans, MD\\
\textbf{機構：}Department of General Internal Medicine, Medical Intensive Care Unit, University Hospitals Leuven, Belgium

%==============================================================
\section{背景知識}
%==============================================================

\subsection{橫膈膜解剖與生理}

\begin{itemize}
    \item \textbf{結構}：薄、圓頂狀肌肉，包含 costal part、crural part 及不收縮的 central tendon
    \item \textbf{收縮機制}：肌纖維活化時，橫膈膜在 zone of apposition 縮短並增厚，使 dome 向下移動
    \item \textbf{功能定義}：產生收縮並生成壓力的能力，具有相當大的儲備能力
\end{itemize}

\subsection{Critical Illness-Associated Diaphragm Weakness}

\begin{itemize}
    \item 在重症病人中常見，甚至在 ICU 早期即可發生
    \item \textbf{危險因子}：Disuse（長期機械通氣）及 inflammation
    \item \textbf{臨床影響}：可能導致 weaning 困難並惡化預後
    \item 約 \textbf{20-30\%} 的機械通氣病人屬於 difficult-to-wean
\end{itemize}

%==============================================================
\section{操作技術}
%==============================================================

\note{
\textbf{通用原則：}
\begin{itemize}
    \item 病人採 \textbf{semi-recumbent position}（30-45°）
    \item \textbf{右側}評估通常較容易且足夠，除非懷疑單側功能異常
    \item 評估時應\textbf{無或僅有最少呼吸器支持}（無法區分主動 vs 被動位移）
    \item 確保有足夠 respiratory drive（如 P0.1 ≥2 cmH$_2$O）
\end{itemize}
}

\subsection{Subcostal Approach}

\begin{longtable}{p{4cm}p{10cm}}
\toprule
\textbf{項目} & \textbf{內容} \\
\midrule
\endfirsthead

\toprule
\textbf{項目} & \textbf{內容} \\
\midrule
\endhead

\bottomrule
\endfoot

\textbf{Probe} & Low frequency 2-5 MHz（abdominal 或 cardiac transducer）\\
\midrule
\textbf{位置} & Mid-clavicular line 下方 subcostally\newline
向頭側傾斜，垂直於 dome \\
\midrule
\textbf{影像特徵} & B-mode：橫膈膜呈 bright line 覆蓋 liver 或 spleen\newline
吸氣時向 probe 移動 \\
\midrule
\textbf{測量方法} & M-mode 測量 tidal breathing 時的 displacement\newline
M-line 需垂直於橫膈膜移動方向 \\
\midrule
\textbf{測量參數} & \textbf{Diaphragm Excursion (DE)}：吸氣時的位移距離 \\
\midrule
\textbf{影像最佳化} & 調整 depth 以最佳捕捉 excursion\newline
調整 gain 以最佳化對比\newline
調整 sweep speed 使至少 3 個呼吸週期在 1 frame 內 \\
\bottomrule
\end{longtable}

\subsection{Intercostal Approach}

\begin{longtable}{p{4cm}p{10cm}}
\toprule
\textbf{項目} & \textbf{內容} \\
\midrule
\endfirsthead

\toprule
\textbf{項目} & \textbf{內容} \\
\midrule
\endhead

\bottomrule
\endfoot

\textbf{Probe} & High frequency 7-12 MHz（linear transducer）\\
\midrule
\textbf{位置} & Antero- 或 mid-axillary line\newline
第 8-11 肋間隙，垂直於胸壁\newline
位於 zone of apposition，盡量與肋間隙平行 \\
\midrule
\textbf{影像特徵} & 三層結構：\newline
• Non-echogenic muscular layer（中間）\newline
• Echogenic pleural lining（上方）\newline
• Echogenic peritoneal lining（下方）\newline
• 中央可見 fibrous layer 呈 bright line \\
\midrule
\textbf{測量方法} & B-mode 或 M-mode 測量 thickness\newline
Calipers 置於 pleural 及 peritoneal lining 內緣（不含在內）\\
\midrule
\textbf{測量參數} & DTee：End-expiration thickness\newline
DTpi：Peak-inspiration thickness\newline
\textbf{DTF = ((DTpi-DTee)/DTee) × 100} \\
\bottomrule
\end{longtable}

%==============================================================
\section{正常值與診斷閾值}
%==============================================================

\subsection{健康人正常值（坐姿、tidal breathing、end-expiration）}

\begin{longtable}{lccc}
\toprule
\textbf{參數} & \textbf{右側男性} & \textbf{右側女性} & \textbf{左側} \\
\midrule
\endfirsthead

\toprule
\textbf{參數} & \textbf{右側男性} & \textbf{右側女性} & \textbf{左側} \\
\midrule
\endhead

\bottomrule
\endfoot

DE & 2.0±0.5 cm & 1.9±0.5 cm & 男 2.2±0.6 cm\newline 女 1.9±0.5 cm \\
\midrule
Thickness & 2.1±0.4 mm & 1.9±0.4 mm & 男 2.0±0.4 mm\newline 女 1.7±0.3 mm \\
\midrule
DTF & 32±15\% & 35±16\% & 男 30±14\%\newline 女 33±15\% \\
\bottomrule
\end{longtable}

\subsection{診斷閾值}

\begin{longtable}{p{5cm}p{9cm}}
\toprule
\textbf{參數} & \textbf{診斷意義} \\
\midrule
\endfirsthead

\toprule
\textbf{參數} & \textbf{診斷意義} \\
\midrule
\endhead

\bottomrule
\endfoot

\textbf{DE <10-15 mm} & Tidal breathing 時 → Diaphragm dysfunction \\
\midrule
\textbf{Paradoxical movement} & 吸氣時向上移動 → Diaphragm paralysis \\
\midrule
\textbf{DTF$_{max}$ <20\%} & Diaphragm dysfunction \\
\midrule
\textbf{DTF <25-33\%} & Predicts weaning failure \\
\midrule
\textbf{DTF <20\%} & Predicts NIV failure \\
\bottomrule
\end{longtable}

%==============================================================
\section{臨床應用}
%==============================================================

\subsection{主要適應症}

\begin{enumerate}
    \item \textbf{Difficult-to-wean 病人的診斷評估}
    \begin{itemize}
        \item SBT 或 extubation failure 的診斷 workup
        \item 識別或排除 diaphragm dysfunction 作為可能原因
        \item 屬於 multimodal diagnostic approach 的一部分
    \end{itemize}

    \item \textbf{Neuromuscular disease 導致的不明原因呼吸衰竭}
    \begin{itemize}
        \item 導致 ICU 入住的呼吸衰竭
        \item 確認 diaphragm weakness
    \end{itemize}

    \item \textbf{Traumatic phrenic nerve injury}
    \begin{itemize}
        \item 確認單側或雙側創傷性膈神經損傷
    \end{itemize}
\end{enumerate}

\subsection{目前的限制}

\note{
\textbf{尚不建議用於以下決策：}
\begin{itemize}
    \item 決定是否執行 SBT 或 extubation
    \item NIV 病人的 intubation 決策
    \item 最佳化 inspiratory effort（「diaphragm protective ventilation」）
\end{itemize}

\textbf{原因：}研究異質性大、多為觀察性、cutoff 值不一致、DTF 與 diaphragm function 相關性不佳
}

\subsection{Difficult-to-Wean 病人的處置建議}

若 diaphragm ultrasound 診斷 diaphragm weakness：

\begin{enumerate}
    \item \textbf{減少 respiratory load}
    \begin{itemize}
        \item 處理 fluid overload
        \item 處理 atelectasis
        \item 引流 pleural effusion
        \item 處理 dynamic hyperinflation
    \end{itemize}

    \item \textbf{若排除 diaphragm weakness}
    \begin{itemize}
        \item 評估 heart failure
        \item 評估 psychological factors
        \item 排除其他 weaning failure 原因
    \end{itemize}
\end{enumerate}

%==============================================================
\section{技術注意事項}
%==============================================================

\subsection{測量品質控制}

\begin{itemize}
    \item 至少取得 \textbf{3 次測量}，差異應 <10\%
    \item DTF 測量需確認 DTpi 和 DTee 來自\textbf{同一呼吸週期}
    \item B-mode 提供較好的空間定位，M-mode 提供較準確的呼吸週期時序
\end{itemize}

\subsection{訓練建議}

\begin{itemize}
    \item 建議至少 \textbf{40 次訓練檢查}才能獨立操作
    \item DT、DTF 及右側橫膈膜 excursion 的測量，在有經驗的操作者手中具有良好的\textbf{可重複性}及\textbf{再現性}
\end{itemize}

%==============================================================
\section{未來發展}
%==============================================================

\begin{enumerate}
    \item \textbf{Asynchrony detection}：偵測機械通氣時的不同步，常被 waveform monitoring 遺漏
    \item \textbf{Operator-independent continuous monitoring}：新型裝置可連續監測橫膈膜
    \item \textbf{Novel ultrasound parameters}：
    \begin{itemize}
        \item Strain 及 stiffness（from speckle tracking）
        \item Shear-wave elastography
        \item 可能成為更可靠的 diaphragm function 指標
    \end{itemize}
\end{enumerate}

%==============================================================
\section{Take-Home Message}
%==============================================================

\begin{enumerate}
    \item \textcolor{red}{\textbf{主要價值：}}非侵入性床邊工具，用於 difficult-to-wean 病人診斷或排除 diaphragm dysfunction
    \item \textcolor{red}{\textbf{首選方法：}}Subcostal approach 測量 DE，DE <10-15 mm 提示 dysfunction
    \item \textcolor{red}{\textbf{其他應用：}}識別 neuromuscular disease 導致的呼吸衰竭、診斷 traumatic phrenic nerve injury
    \item \textcolor{red}{\textbf{注意限制：}}目前不建議用於 SBT/extubation 或 intubation 決策
\end{enumerate}

%==============================================================
\section{縮寫對照表}
%==============================================================

\begin{longtable}{p{3cm}p{11cm}}
\toprule
\textbf{縮寫} & \textbf{全名} \\
\midrule
\endfirsthead

\toprule
\textbf{縮寫} & \textbf{全名} \\
\midrule
\endhead

\bottomrule
\endfoot

DE & Diaphragm Excursion (橫膈膜位移) \\
DTee & Diaphragm Thickness at End-Expiration \\
DTpi & Diaphragm Thickness at Peak-Inspiration \\
DTF & Diaphragm Thickening Fraction (橫膈膜增厚分數) \\
SBT & Spontaneous Breathing Trial (自發呼吸試驗) \\
NIV & Non-Invasive Ventilation (非侵入性通氣) \\
P0.1 & Airway Occlusion Pressure at 0.1 second \\
\bottomrule
\end{longtable}

%==============================================================
\section{參考文獻}
%==============================================================

\begin{enumerate}
    \item Hermans G, Demoule A, Heunks L. How I perform diaphragmatic ultrasound in the intensive care unit. \href{https://doi.org/10.1007/s00134-024-07688-x}{\textit{Intensive Care Med}. 2024;50:2175-2178.}

    \item Haaksma ME, Smit JM, Boussuges A, et al. EXpert consensus On Diaphragm UltraSonography in the critically ill (EXODUS). \href{https://doi.org/10.1186/s13054-022-03977-5}{\textit{Crit Care}. 2022;26:99.}

    \item Tuinman PR, Jonkman AH, Dres M, et al. Respiratory muscle ultrasonography: methodology, basic and advanced principles and clinical applications in ICU and ED patients. \href{https://doi.org/10.1007/s00134-020-05892-7}{\textit{Intensive Care Med}. 2020;46:594-605.}

    \item Dres M, Goligher EC, Heunks LMA, Brochard LJ. Critical illness-associated diaphragm weakness. \href{https://doi.org/10.1007/s00134-017-4928-4}{\textit{Intensive Care Med}. 2017;43:1441-1452.}
\end{enumerate}

\end{document}
